\chapter*{INTRODUCCIÓN}\addcontentsline{toc}{chapter}{INTRODUCCIÓN}\label{introduccion}

La congestión vehicular en Lima Metropolitana es uno de los problemas urbanos más críticos que afecta
diariamente a millones de personas. Un informe publicado por El Comercio en 2024 ubica a Lima entre las
ciudades con mayor nivel de congestión vehicular en América Latina, con una velocidad promedio cercana a
14.5 km/h. Este diagnóstico es consistente con los análisis de movilidad urbana de organismos
institucionales, como el Banco Mundial (2024). Este problema no solo representa una pérdida significativa de tiempo para los conductores, sino
que también genera un impacto económico relevante asociado a horas-hombre improductivas, mayor consumo de
combustible y menor eficiencia logística. Por otra parte, el transporte contribuye de manera importante a
las emisiones contaminantes metropolitanas (Banco Mundial, 2024), por lo que la mejora del control del
tránsito también tiene implicancias ambientales.

Una de las principales causas de esta congestión es la ineficiencia del sistema semafórico actual, que en
gran medida opera con ciclos fijos. En Lima Metropolitana existen cerca de 1\,200 intersecciones
semaforizadas bajo administración municipal, de las cuales una amplia mayoría opera de forma descoordinada
y sin gestión centralizada (Diario Correo, 2020). La mayor parte de los semáforos no tiene capacidad de
adaptarse a las variaciones del flujo vehicular en tiempo real; solo un conjunto reducido de intersecciones
cuenta con semáforos denominados «inteligentes», mientras que el resto opera de manera aislada y sin
conexión a un centro de control. Este modelo rígido aumenta los tiempos de espera, da lugar a una
utilización ineficiente de recursos como el combustible y el tiempo de los conductores, e incrementa los
niveles de contaminación (Sacyr, s.f.).

Frente a este escenario, la modernización de los sistemas semafóricos mediante tecnologías inteligentes es
una necesidad para optimizar la movilidad urbana y reducir sus impactos negativos. Sin embargo, la
sustitución completa de la infraestructura instalada resulta costosa y de implementación lenta. La presente
tesis aborda el problema como un desafío de integración mecatrónica: el sistema debe observar el estado de
la intersección, transformar mediciones en variables de tráfico, calcular ajustes dentro de límites
operativos, conservar capacidad de decisión local ante pérdida de conectividad y comunicar telemetría sin
comprometer la seguridad funcional del controlador vial. Para ello se diseña un módulo IoT complementario
con visión por computadora, procesamiento en el borde, control difuso y una plataforma en la nube destinada
a supervisión, almacenamiento histórico y distribución controlada de parámetros.

El diseño se sustenta en las lecciones aprendidas de países con sistemas semafóricos avanzados y de
experiencias exitosas en Latinoamérica, adaptando las mejores prácticas a las necesidades locales. El
documento se organiza en siete capítulos. El primero presenta los antecedentes, la formulación del
problema, los objetivos, el alcance y la metodología de diseño (VDI 2221 y VDI 2206) que guía el trabajo.
El segundo revisa el estado del arte en sensado y detección vehicular, control adaptativo basado en lógica
difusa, arquitecturas edge--nube, ciberseguridad y validación por simulación. El tercero desarrolla el
diseño conceptual a partir de las necesidades de los actores, la lista de requerimientos, la estructura de
funciones y la matriz morfológica, hasta seleccionar la solución óptima. El cuarto y el quinto detallan,
respectivamente, el diseño electrónico ---selección de sensores, actuadores y el enlace de comunicación con
la nube--- y el diseño mecánico ---materiales, análisis estructural y el mecanismo de orientación de la
cámara---. El sexto define la etapa de control, organizada en tres capas, e introduce el Índice de
Congestión Vehicular (ICV) y el controlador difuso que ajusta los tiempos de fase. Finalmente, el séptimo
integra los subsistemas y valida la propuesta mediante simulaciones en SUMO sobre escenarios de Lima,
comparando el control difuso con el de tiempo fijo y con líneas base avanzadas, evaluando el
comportamiento ante pérdida de conectividad y presentando el análisis de costos. El documento cierra con
las conclusiones y recomendaciones derivadas del trabajo.
